\chapter{Frontend — interface StockFlow}
\label{ch:frontend}

\section{Stack technique}

L'interface utilisateur est une \textbf{single-page application} construite avec :
\begin{itemize}
  \item \textbf{React 18} et \textbf{TypeScript 5.6} ;
  \item \textbf{Vite 6} comme bundler et serveur de développement ;
  \item \textbf{React Router 7} pour le routage côté client ;
  \item \textbf{Tailwind CSS 4} via le plugin \texttt{@tailwindcss/vite} ;
  \item \textbf{lucide-react} pour l'iconographie ;
  \item \textbf{Vitest} et React Testing Library pour les tests.
\end{itemize}

Le nom produit affiché est \stockflow{} : thème sombre, accents cyan, typographie Outfit (chargée depuis Google Fonts dans \texttt{index.html}).

\section{Architecture frontale}

\subsection{Organisation des sources}

\begin{table}[htbp]
  \centering
  \caption{Structure \texttt{frontend/src/}}
  \label{tab:frontend-structure}
  \begin{tabular}{@{}lp{9cm}@{}}
    \toprule
    \textbf{Dossier} & \textbf{Contenu} \\
    \midrule
    \texttt{pages/} & Une page par route métier \\
    \texttt{composants/Layout.tsx} & Shell navigation + zone \texttt{Outlet} \\
    \texttt{composants/ui/} & Bibliothèque de composants réutilisables \\
    \texttt{contexte/} & Auth et toasts globaux \\
    \texttt{hooks/useFetch.ts} & Chargement données API avec état \\
    \texttt{services/api.ts} & Client HTTP, gestion du jeton JWT \\
    \bottomrule
  \end{tabular}
\end{table}

\subsection{Routage et protection}

\texttt{App.tsx} définit :
\begin{itemize}
  \item \texttt{/connexion} — publique ;
  \item routes imbriquées sous \texttt{Layout} — protégées par \texttt{RoutesProtegees} qui lit \texttt{AuthContext} ;
  \item redirection vers \texttt{/connexion} si non authentifié ;
  \item affichage d'un chargeur pendant la vérification du jeton (\texttt{GET /auth/moi}).
\end{itemize}

\section{Client API}

\texttt{services/api.ts} centralise les appels :
\begin{itemize}
  \item base URL depuis \texttt{import.meta.env.VITE\_API\_URL} (injectée au build Docker) ;
  \item jeton stocké dans \texttt{sessionStorage} sous la clé \texttt{jeton\_acces} ;
  \item fonction \texttt{appelApi<T>()} ajoutant l'en-tête \texttt{Authorization} et gérant les erreurs 401 (déconnexion implicite).
\end{itemize}

En mode cluster, le frontend est servi depuis \texttt{http://localhost} et les appels partent vers \texttt{http://localhost/api/...} via l'Ingress — d'où le build avec \texttt{VITE\_API\_URL=http://localhost}.

\section{Pages fonctionnelles}

\subsection{Connexion (\texttt{PageConnexion.tsx})}

Formulaire email / mot de passe, valeurs préremplies en démo. Appel \texttt{POST /auth/connexion}, stockage du jeton, toast de succès, navigation vers le tableau de bord. Redirection automatique si session déjà valide.

\subsection{Tableau de bord}

Consomme \texttt{GET /api/v1/tableau-de-bord}. Affiche quatre cartes statistiques (produits, catégories, mouvements, alertes) et une liste des produits sous seuil avec barres de progression visuelles.

\subsection{Produits}

Liste filtrable (nom, SKU, catégorie), tableau avec badge d'alerte si stock bas, modal de création (POST), suppression avec confirmation (DELETE, visible au survol).

\subsection{Catégories}

Grille de cartes avec description tronquée, modal de création, suppression au survol.

\subsection{Mouvements}

Timeline verticale : entrées (vert) et sorties (rouge), date et motif. Modal pour enregistrer un mouvement (type, produit, quantité, motif optionnel).

\subsection{Alertes}

Cartes par produit en alerte avec jauge et pourcentage par rapport au seuil ; lien vers les mouvements pour corriger le stock.

\section{Composants UI}

La couche \texttt{composants/ui/} homogénéise l'interface :
\begin{itemize}
  \item \texttt{Bouton} — variantes primaire, secondaire, danger, fantôme ; état chargement ;
  \item \texttt{Modal} — fermeture Escape, tailles md/lg ;
  \item \texttt{Champ}, \texttt{Input}, \texttt{Select}, \texttt{Textarea} ;
  \item \texttt{CarteStat}, \texttt{Badge}, \texttt{Chargeur}, \texttt{EtatVide}, \texttt{BarreRecherche}.
\end{itemize}

\section{État global}

\texttt{AuthContext} expose \texttt{utilisateur}, \texttt{chargement}, \texttt{connexion}, \texttt{deconnexion}. Au montage, si un jeton existe en session, appel \texttt{/auth/moi} pour hydrater le profil.

\texttt{ToastContext} permet d'afficher des notifications succès / erreur / info en bas à droite (auto-dismiss après 4 secondes).

\section{Layout responsive}

\texttt{Layout.tsx} combine :
\begin{itemize}
  \item barre latérale fixe sur grand écran (navigation, profil, déconnexion) ;
  \item menu hamburger et tiroir sur mobile ;
  \item animation \texttt{animate-fade-in} à chaque changement de route (\texttt{key=\{location.pathname\}}).
\end{itemize}

\section{Build et livraison}

Le \texttt{frontend/Dockerfile} est multi-stage :
\begin{enumerate}
  \item Stage \texttt{build} : \texttt{node:20-alpine}, \texttt{npm install}, \texttt{npm run build} (TypeScript + Vite).
  \item Stage runtime : \texttt{nginx:alpine}, copie de \texttt{dist/} et \texttt{nginx.conf}.
\end{enumerate}

\texttt{nginx.conf} configure le fallback SPA (\texttt{try\_files}) et, en Compose, un proxy \texttt{/api/} vers le service \api{} — en cluster, c'est l'Ingress qui assume ce routage.

\section{Parcours utilisateur et états d'interface}

Chaque page métier suit le même patron : chargement (\texttt{Chargeur}), état vide (\texttt{EtatVide}) si la liste est vide, puis contenu principal. Les actions destructives (suppression) utilisent une confirmation implicite via bouton dédié et toast d'erreur en cas d'échec API.

\begin{table}[htbp]
  \centering
  \caption{Pages et endpoints consommés}
  \label{tab:frontend-pages}
  \small
  \begin{tabular}{@{}ll@{}}
    \toprule
    \textbf{Route} & \textbf{Endpoints principaux} \\
    \midrule
    \texttt{/connexion} & \texttt{POST /auth/connexion} \\
    \texttt{/} (dashboard) & \texttt{GET /tableau-de-bord} \\
    \texttt{/produits} & \texttt{GET/POST/DELETE /produits} \\
    \texttt{/categories} & \texttt{GET/POST/DELETE /categories} \\
    \texttt{/mouvements} & \texttt{GET/POST /mouvements} \\
    \texttt{/alertes} & \texttt{GET /produits/alertes} \\
    \bottomrule
  \end{tabular}
\end{table}

Le hook \texttt{useFetch} encapsule le cycle \emph{chargement → données → erreur}, ce qui évite de dupliquer les \texttt{useEffect} sur chaque page.

\section{Tests frontend}

\texttt{App.test.tsx} vérifie le rendu du titre « Connexion » avec les fournisseurs de contexte et un \texttt{MemoryRouter}. La CI exécute \texttt{npm run lint} (ESLint) et \texttt{npm test -- --run} (Vitest).

Les règles ESLint (config flat dans \texttt{eslint.config.js}) couvrent les hooks React et les imports inutilisés ; elles complètent TypeScript qui reste la première barrière de typage au build \texttt{tsc -b}.
